\documentclass[12pt,a4paper]{article}
\usepackage[T2A]{fontenc}
\usepackage[utf8]{inputenc}
\usepackage[russian]{babel}
\usepackage{geometry}
\geometry{left=21mm,right=21mm,top=23mm,bottom=24mm}
\usepackage{amsmath,amssymb,mathtools}
\usepackage{array,booktabs,longtable}
\usepackage{xcolor}
\usepackage{tikz}
\usepackage[unicode,colorlinks=true,linkcolor=blue!55!black]{hyperref}
\setlength{\parindent}{0pt}
\setlength{\parskip}{6pt}
\renewcommand{\arraystretch}{1.28}
\newcommand{\err}{\textcolor{red}{\textbf{Ошибка:}}}
\newcommand{\idea}{\textcolor{blue!60!black}{\textbf{Ключевая идея:}}}
\newcommand{\result}{\textcolor{teal!60!black}{\textbf{Итог:}}}
\newcommand{\reviewhead}[1]{\medskip\textbf{#1}\par}
\begin{document}
\thispagestyle{empty}
\vspace*{0.27\textheight}
\begin{center}
{\Huge\bfseries Ревью домашней работы}\par
\vspace{1.7cm}
{\large Автор: Лёвин Артём Александрович}\par
\vspace{0.7cm}
{\large 19 октября 2026 года}
\end{center}
\vfill
\begin{center}
\begin{tikzpicture}
\draw[blue!45!black,line width=0.85pt] (-3.3,0) .. controls (-1.7,0.52) and (-0.4,-0.5) .. (0.6,0) .. controls (1.6,0.48) and (2.45,-0.2) .. (3.3,0.05);
\end{tikzpicture}
\end{center}
\newpage
\section*{Сводная таблица ошибок}
\small
\begin{longtable}{@{}>{\centering\arraybackslash}p{0.075\linewidth}p{0.17\linewidth}p{0.26\linewidth}p{0.18\linewidth}p{0.18\linewidth}@{}}
\toprule
\textbf{№} & \textbf{Тема} & \textbf{Суть ошибки} & \textbf{Ваш ответ} & \textbf{Правильный ответ}\\
\midrule
\endfirsthead
\toprule
\textbf{№} & \textbf{Тема} & \textbf{Суть ошибки} & \textbf{Ваш ответ} & \textbf{Правильный ответ}\\
\midrule
\endhead
\hyperref[task:2]{2} & Оценка дроби & Неправильно упорядочены значения обратной зависимости. & $6<\frac{24}{t}<4$ & $4<\frac{24}{t}<6$ \\
\hyperref[task:4]{4} & Границы произведения и частного & Указаны границы произведения, но не исследовано частное $k/n$. & $-6\le kn\le3$ & Для $k/n$ конечных границ нет. \\
\hyperref[task:8a]{8а} & Простейшее неравенство & Вместо плюса в условии переписан минус. & $x<x-6$ & $x\in\mathbb R$ \\
\hyperref[task:9]{9} & Линейное неравенство, целые числа & Получен промежуток, не перечислены отрицательные целые решения. & $x>-\frac83$ & $x=-2,-1$ \\
\bottomrule
\end{longtable}
\normalsize
\newpage
\thispagestyle{empty}
\vspace*{\fill}
\begin{center}
{\LARGE\bfseries Разбор ошибок}
\end{center}
\vspace*{\fill}

\newpage
\phantomsection\label{task:2}
\section*{Задание № 2. Оценка дроби}
\reviewhead{Текст задания}
Известно, что $4<t<6$. Укажите промежуток для $\dfrac{24}{t}$.

\reviewhead{Ваше решение}
В работе записано:
\[
4<t<6,\qquad
\frac{24}{4}<\frac{24}{t}<\frac{24}{6},
\]
а затем
\[
6<\frac{24}{t}<4.
\]

\reviewhead{Ваш ответ}
$\displaystyle 6<\frac{24}{t}<4$.

\reviewhead{Ошибка}
\err{} Нарушен порядок границ при переходе от положительного числа $t$ к дроби $24/t$.

\reviewhead{Где именно возникает ошибка}
Последняя правильная запись: $4<t<6$. Первый неверный переход:
\[
\frac{24}{4}<\frac{24}{t}<\frac{24}{6}.
\]
Знак «меньше» между крайними дробями неверен, потому что $24/4=6$, а $24/6=4$: число шесть не меньше четырёх.

\reviewhead{Почему это неверно}
Числитель дроби остаётся одним и тем же положительным числом, а знаменатель увеличивается. При этом значение дроби \emph{уменьшается}. Например, $24/4=6$, $24/5=4{,}8$, $24/6=4$. Следовательно, для дробей с постоянным положительным числителем значения идут в обратном порядке по отношению к положительным знаменателям.

\idea{} Если положительный знаменатель увеличивается, дробь с фиксированным положительным числителем уменьшается.

\reviewhead{Исправление от последнего правильного шага}
Возвращаемся к исходному двойному неравенству $4<t<6$. Поскольку все три числа положительны, при переходе к обратным величинам оба знака меняются по направлению:
\[
\frac14>\frac1t>\frac16.
\]
Умножение каждой части на положительное число $24$ не меняет направления знаков:
\[
\frac{24}{4}>\frac{24}{t}>\frac{24}{6}.
\]
После вычисления получаем $6>\dfrac{24}{t}>4$.

\reviewhead{Полное правильное решение}
\begin{align*}
4<t<6
&\quad\Longrightarrow\quad \frac16<\frac1t<\frac14,\\
&\quad\Longrightarrow\quad \frac{24}{6}<\frac{24}{t}<\frac{24}{4},\\
&\quad\Longrightarrow\quad \boxed{4<\frac{24}{t}<6}.
\end{align*}
Крайние значения не включаются: по условию $t$ строго больше четырёх и строго меньше шести.

\reviewhead{Проверка}
Возьмём допустимое значение $t=5$. Тогда $24/t=4{,}8$: это число находится между четырьмя и шестью. Для Вашей записи пришлось бы одновременно требовать, чтобы $4{,}8$ было больше шести и меньше четырёх, что невозможно.

\result{} $\displaystyle \frac{24}{t}\in(4;6)$.

\reviewhead{Задача на закрепление}
Известно, что $5<s<8$. Укажите промежуток для $\dfrac{40}{s}$.

\newpage
\phantomsection\label{task:4}
\section*{Задание № 4. Границы произведения и частного}
\reviewhead{Текст задания}
При условиях $-2\le k\le1$ и $-1\le n\le3$ найдите границы $kn$. Можно ли указать конечные нижнюю и верхнюю границы $k/n$ для всех допустимых $k,n$? Объясните.

\reviewhead{Ваше решение}
Вы рассмотрели произведение и вычислили значения на сочетаниях концов промежутков:
\[
(-2)(-1)=2,\quad (-2)\cdot3=-6,
\]
\[
1\cdot(-1)=-1,\quad 1\cdot3=3.
\]
После этого указаны границы $-6$ и $3$. Исследование дроби $k/n$ в решении отсутствует.

\reviewhead{Ваш ответ}
Для произведения указано $-6\le kn\le3$. Ответ о частном $k/n$ не записан.

\reviewhead{Ошибка}
\err{} Задание выполнено не полностью: не проверено, определено ли частное для всех допустимых значений, и существуют ли у него конечные границы.

\reviewhead{Где именно возникает ошибка}
Последний завершённый шаг --- получение границ произведения. После него необходимо перейти ко второй части вопроса, однако вывод о дроби $k/n$ не сделан.

\reviewhead{Почему это неверно}
Границы произведения нельзя автоматически переносить на частное: произведение $kn$ и дробь $k/n$ ведут себя по-разному. В частности, произведение двух ограниченных чисел остаётся ограниченным. Но для дроби значение знаменателя может стать сколь угодно близким к нулю, из-за чего модуль дроби растёт без ограничения. Более того, при $n=0$ деление вообще не определено. Число ноль разрешено условием $-1\le n\le3$, и это нужно учитывать отдельно.

\idea{} Сначала проверить знаменатель на возможность обращения в ноль. Если ненулевой знаменатель может сколь угодно близко подходить к нулю, а числитель остаётся ненулевым, частное может не иметь конечных границ.

\reviewhead{Исправление от последнего правильного шага}
Сохраняем полученные границы $-6\le kn\le3$. Для частного выбираем разрешённое значение $k=1$. Положим $n=1/m$, где $m$ --- любое натуральное число. Тогда $n\in[-1;3]$, $n\ne0$, причём
\[
\frac{k}{n}=\frac{1}{1/m}=m.
\]
При увеличении $m$ значения частного сколь угодно велики; следовательно, конечной верхней границы нет.

Для нижней границы возьмём $k=1$, $n=-1/m$. Эти значения также допустимы. Тогда
\[
\frac{k}{n}=\frac{1}{-1/m}=-m.
\]
Можно получить сколь угодно большие по модулю отрицательные значения. Значит, конечной нижней границы тоже нет.

\reviewhead{Полное правильное решение}
\textbf{Шаг 1. Границы произведения.} При фиксированном $n$ выражение $kn$ линейно по $k$, а при фиксированном $k$ --- линейно по $n$. Поэтому на прямоугольнике допустимых значений наименьшее и наибольшее значения произведения достигаются при сочетаниях концов промежутков. Перебираем четыре случая:
\[
\begin{array}{c|cc}
 & n=-1 & n=3\\ \hline
k=-2 & 2 & -6\\
k=1 & -1 & 3
\end{array}
\]
Самое маленькое значение равно $-6$ (при $k=-2$, $n=3$), самое большое равно $3$ (при $k=1$, $n=3$). Поэтому
\[
\boxed{-6\le kn\le3}.
\]
\textbf{Шаг 2. Область определения частного.} Дробь $k/n$ имеет смысл только при $n\ne0$. В исходном промежутке для $n$ ноль содержится, поэтому говорить о значении частного при \emph{всех} исходно разрешённых парах без дополнительного ограничения нельзя.

\textbf{Шаг 3. Проверка верхней и нижней границы.} Даже если оставить только пары с $n\ne0$, можно брать $k=1$ и приближать $n$ к нулю справа или слева. В первом случае $k/n$ становится сколь угодно большим положительным числом; во втором --- сколь угодно большим по модулю отрицательным числом. Значит, ни конечной верхней, ни конечной нижней границы не существует.

\reviewhead{Проверка}
Проверяем примерами с допустимыми значениями: при $k=1$, $n=0{,}01$ частное равно $100$, а при $k=1$, $n=-0{,}01$ оно равно $-100$. Если вместо сотни требуется число большей величины, достаточно выбрать знаменатель ещё ближе к нулю. Никакое фиксированное конечное число не ограничит все такие частные.

\result{} $-6\le kn\le3$; дробь $k/n$ при $n=0$ не определена, а среди допустимых пар с $n\ne0$ не имеет ни конечной нижней, ни конечной верхней границы.

\reviewhead{Задача на закрепление}
Пусть $-3\le p\le2$ и $-2\le q\le4$. Найдите границы $pq$. Можно ли указать конечные нижнюю и верхнюю границы $p/q$ для всех допустимых значений, при которых частное определено? Объясните.

\newpage
\phantomsection\label{task:8a}
\section*{Задание № 8а. Множество решений неравенства}
\reviewhead{Текст задания}
Укажите множество решений неравенства $x+6>x$.

\reviewhead{Ваше решение}
В первой строке работы вместо заданного выражения записано $x-6>x$. Далее выполнен переход к $x<x-6$. Отдельное множество решений для пункта «а» не указано.

\reviewhead{Ваш ответ}
Записано промежуточное неравенство $x<x-6$; итоговое множество решений не приведено.

\reviewhead{Ошибка}
\err{} При переписывании условия знак перед шестёркой заменён: вместо сложения получилось вычитание. Изменилось само неравенство, которое требуется решить.

\reviewhead{Где именно возникает ошибка}
Последняя достоверно правильная математическая запись --- исходное условие $x+6>x$. Первый неверный шаг --- замена этой записи на $x-6>x$ без равносильного преобразования.

\reviewhead{Почему это неверно}
Выражения $x+6$ и $x-6$ отличаются на двенадцать единиц. Следовательно, нельзя заменить одно другим, не изменяя задачу. Более того, исходное неравенство выполняется при любом $x$, поскольку прибавление положительного числа шесть делает число больше. А неравенство $x-6>x$, наоборот, не выполняется ни при каком $x$: вычитание шести делает число меньше. Это два противоположных случая.

\idea{} Перед преобразованиями сохранять знаки в условии. При вычитании одного и того же выражения из обеих частей знак неравенства не меняется.

\reviewhead{Исправление от последнего правильного шага}
Возвращаемся к $x+6>x$. Вычитаем $x$ из левой и правой частей:
\[
x+6-x>x-x.
\]
Остаётся $6>0$. Это верное числовое неравенство, которое не зависит от значения $x$.

\reviewhead{Полное правильное решение}
\begin{align*}
x+6&>x,\\
x+6-x&>x-x,\\
6&>0.
\end{align*}
Последнее утверждение всегда верно. Поэтому исходное неравенство не накладывает на $x$ никаких ограничений: подходит любое действительное число.

\reviewhead{Проверка}
При $x=-10$ получаем $-4>-10$; при $x=0$ получаем $6>0$; при $x=100$ получаем $106>100$. Во всех случаях прибавление шести увеличивает исходное число. Общий переход к $6>0$ доказывает это для всех действительных $x$.

\result{} Множество решений: $\boxed{x\in\mathbb R}$, или $(-\infty;+\infty)$.

\reviewhead{Задача на закрепление}
Укажите множество решений неравенства $y+9>y$.

\newpage
\phantomsection\label{task:9}
\section*{Задание № 9. Отрицательные целые решения}
\reviewhead{Текст задания}
Найдите все отрицательные целые решения неравенства
\[
\frac{13x-1}{15}-\frac{2x-1}{5}<x-\frac{x-2}{3}.
\]

\reviewhead{Ваше решение}
Вы правильно привели дроби к общему знаменателю и получили:
\begin{align*}
(13x-1)-(6x-3)&<15x-(5x-10),\\
7x+2&<10x+10,\\
-3x&<8,\\
x&>-\frac83.
\end{align*}
На этом решение заканчивается.

\reviewhead{Ваш ответ}
$\displaystyle x>-\frac83$.

\reviewhead{Ошибка}
\err{} Получен верный промежуток решений неравенства, но не выполнено дополнительное требование: найти \emph{все отрицательные целые} решения.

\reviewhead{Где именно возникает ошибка}
Последний правильный шаг --- вывод $x>-8/3$. После него необходимо отобрать числа, которые одновременно удовлетворяют этому неравенству, являются целыми и меньше нуля. Эта часть решения пропущена.

\reviewhead{Почему это неверно}
Промежуток $x>-8/3$ содержит не только отрицательные целые числа, но и дробные, ноль и положительные значения. Например, число $1/2$ ему удовлетворяет, но не является отрицательным целым числом. Поэтому запись одного промежутка не отвечает на вопрос задания. Требуется пересечение множества решений неравенства с множеством отрицательных целых чисел.

\idea{} После решения неравенства отдельно применить оба ограничения из текста: число должно быть целым и отрицательным.

\reviewhead{Исправление от последнего правильного шага}
Из результата $x>-8/3$ имеем
\[
-\frac83=-2\frac23.
\]
Отрицательные целые числа, расположенные правее числа минус две целых две трети, --- это $-2$ и $-1$. Число $-3$ уже левее найденной границы, число $0$ не является отрицательным.

\reviewhead{Полное правильное решение}
\textbf{Шаг 1.} Приведём обе части к знаменателю пятнадцать:
\[
\frac{13x-1}{15}-\frac{3(2x-1)}{15}
<\frac{15x}{15}-\frac{5(x-2)}{15}.
\]
\textbf{Шаг 2.} Умножим обе части на положительное число пятнадцать. Направление знака не меняется:
\[
13x-1-3(2x-1)<15x-5(x-2).
\]
\textbf{Шаг 3.} Раскроем скобки, учитывая минусы перед ними:
\begin{align*}
13x-1-6x+3&<15x-5x+10,\\
7x+2&<10x+10.
\end{align*}
\textbf{Шаг 4.} Перенесём слагаемые с $x$ влево, числа вправо:
\begin{align*}
7x-10x&<10-2,\\
-3x&<8.
\end{align*}
\textbf{Шаг 5.} Разделим обе части на отрицательное число $-3$, поэтому знак строго меняет направление:
\[
x>-\frac83.
\]
\textbf{Шаг 6.} Учитываем условие задачи: $x$ --- отрицательное целое число. Подходят только два значения:
\[
\boxed{x\in\{-2;-1\}}.
\]

\reviewhead{Проверка}
Упрощённое исходное неравенство $-3x<8$ ему равносильно. При $x=-2$ получаем $6<8$ --- верно. При $x=-1$ получаем $3<8$ --- верно. Для ближайшего меньшего целого числа $x=-3$ имеем $9<8$ --- неверно. При дальнейшем уменьшении отрицательного целого $x$ выражение $-3x$ только увеличивается, поэтому других отрицательных целых решений нет.

\result{} Все отрицательные целые решения: $\boxed{-2\text{ и }-1}$.

\reviewhead{Задача на закрепление}
Найдите все отрицательные целые решения неравенства
\[
\frac{9x-2}{12}-\frac{x-1}{4}<x-\frac{x-2}{3}.
\]

\vfill
\begin{center}
\footnotesize Лёвин Артём Александрович эксклюзивно для Екатерины Скелли
\end{center}
\end{document}
